\documentclass[12pt]{article}
\usepackage{lmodern}
\usepackage{geometry}
\usepackage{enumitem}
\setlist[description]{leftmargin=1cm,labelindent=1cm}
\usepackage{amsmath}
\usepackage{amsthm}
\usepackage{amssymb}
\usepackage{mathtools}
\usepackage[portuguese]{babel}

\theoremstyle{plain}
\newtheorem{thm}{Teorema}[section]
\newtheorem{prop}[thm]{Proposição}
\newtheorem{lem}[thm]{Lema}
\newtheorem{cor}[thm]{Corolário}

\theoremstyle{definition}
\newtheorem{defn}[thm]{Definição}
\newtheorem{exmp}[thm]{Exemplo}
\newtheorem{prblm}[thm]{Problema}
\newtheorem{prcp}[thm]{Princípio}

\theoremstyle{remark}
\newtheorem*{rem}{Nota}

\geometry{
 a4paper,
 left=3cm,
 top=2cm,
 right=2cm,
 bottom=2cm
}

\begin{document}
\begin{center}
    \textbf{\large UNIVERSIDADE FEDERAL DO RIO GRANDE DO SUL} \\
    \textbf{\large INSTITUTO DE MATEMÁTICA E ESTATÍSTICA} \\
    \vspace{3cm}
    \textbf{\Large MATEMÁTICA DISCRETA I} \\
    \vspace{1cm}
    \textbf{\large LISTA DE RESULTADOS} \\
    \vspace{3cm}

    \textbf{LEONARDO SUAREZ}

    \vfill
    \textbf{\large Porto Alegre, RS} \\
    \textbf{\large \today}
\end{center}

\newpage
\tableofcontents
\newpage

\section{MÉTODOS DE CONTAGEM}
\begin{prcp} (princípio da adição)
    Se existem \(r_1\) objetos distintos no primeiro conjunto, \(r_2\) no segundo conjunto, \ldots, \(r_n\) no n-ésimo conjunto e esses conjuntos são \emph{disjuntos dois a dois}, então o número de maneiras de selecionar um objeto de um dos \(n\) conjuntos é \(r_1 + r_2 + \cdots + r_n\).
\end{prcp}

\begin{prcp} (princípio da multiplicação)
    Se um processo pode ser dividido em \(n\) etapas, onde a primeira etapa pode ser realizada de \(r_1\) maneiras distintas, a segunda etapa pode ser realizada de \(r_2\) maneiras distintas, \ldots, e a n-ésima etapa pode ser realizada de \(r_n\) maneiras distintas, então o número total de maneiras de realizar o processo é \(r_1 \cdot r_2 \cdots r_n\), dado que a escolha de cada etapa é independente das outras.
\end{prcp}

\begin{prcp} (princípio da composição de conjuntos)
    Dado qualquer conjunto \(S\) construído pelo princípio multiplicativo, é preciso saber dizer quais elementos de \(S\) são da primeira parte de \(S\) e quais são da segunda parte.
\end{prcp}

\begin{thm} (arranjos com repetições)
    Se existem \(n\) objetos, com \(r_1\) do tipo 1, \(r_2\) do tipo 2, \ldots, \(r_m\) do tipo \(m\), onde \(r_1 + r_2 + \cdots + r_m = n\), então o número de arranjos distintos que podem ser formados com esses \(n\) objetos é dado por
    \begin{align}
        P(n;r_1,r_2,\ldots,r_m) &= \binom{n}{r_1}\binom{n-r_1}{r_2}\cdots\binom{n-r_1-r_2-\cdots-r_{m-1}}{r_m} \\
        &= \frac{n!}{r_1!r_2!\cdots r_m!}.
    \end{align}
\end{thm}

\begin{thm} (seleções com repetições)
    O número de seleções com repetição de \(r\) objetos dentre \(n\) tipos distintos é \(C(r+n-1,r)\).
\end{thm}

\begin{prop} (formas equivalentes para seleções com repetições)
    \begin{enumerate}
        \item O número de maneiras de selecionar \(r\) objetos com repetição dentro \(n\) tipos distintos.
        \item O número de maneiras de distribuir \(r\) objetos idênticos em \(n\) caixas distintas.
        \item O número de soluções inteiras não negativas da equação
        \begin{equation}
            x_1 + x_2 + \cdots + x_n = r.
        \end{equation}
    \end{enumerate}
\end{prop}

\begin{thm} (teorema binomial)
    Para quaisquer inteiros não-negativos \(n\),
    \begin{equation}
        (x+y)^n = \sum_{k=0}^{n} \binom{n}{k} x^k y^{n-k}.
    \end{equation}
\end{thm}

\begin{prop} (identidades binomiais)
    \begin{gather}
        \binom{n}{0}+\binom{n}{1}+\binom{n}{2}+\cdots+\binom{n}{n} = 2^n \\
        \binom{n}{0}+\binom{n+1}{1}+\binom{n+2}{2}+\cdots+\binom{n+r}{r} = \binom{n+r+1}{r+1} \\
        \binom{r}{r}+\binom{r+1}{r}+\binom{r+2}{r}+\cdots+\binom{n}{r} = \binom{n+1}{r+1} \\
        \binom{n}{0}^2 + \binom{n}{1}^2 + \binom{n}{2}^2 + \cdots + \binom{n}{n}^2 = \binom{2n}{n} \\
        \sum_{k=0}^{r} \binom{m}{k}\binom{n}{r-k} = \binom{m+n}{r} \\
        \sum_{k=0}^{m} \binom{m}{k}\binom{n}{r+k} = \binom{m+n}{m+r} \\
        \sum_{k=n-s}^{m-r} \binom{m-k}{r}\binom{n+k}{s} = \binom{m+n+1}{r+s+1} \\
    \end{gather}
\end{prop}


\section{INCLUSÃO-EXCLUSÃO}

\begin{thm} (princípio da inclusão-exclusão)
    Sejam \(A_1, A_2, \ldots, A_n\) subconjuntos finitos de um conjunto universal \(U\). Seja \(S_k\) a soma de todas as interseções de \(k\) conjuntos dentre os \(A_i\). então
    \begin{equation}
        \left| \bigcup_{i=1}^{n} A_i \right| = S_1 - S_2 + S_3 - S_4 + \cdots + (-1)^{n+1} S_n.
    \end{equation}
\end{thm}

\section{FUNÇÕES GERADORAS}

\section{RELAÇÕES DE RECORRÊNCIA}

\section{TEORIA DE GRAFOS}

\end{document}